\documentclass[12pt,a4paper]{article}

% ------------------------------------------------
% Pacotes
% ------------------------------------------------
\usepackage[utf8]{inputenc}
\usepackage[T1]{fontenc}
\usepackage[brazil]{babel}
\usepackage{graphicx}
\usepackage[table]{xcolor}
\usepackage{array}
\usepackage{tabularx}
\usepackage{geometry}
\usepackage{fancyhdr}
\usepackage{mathptmx} % Times New Roman
\usepackage{courier} % Adiciona suporte a negrito na fonte monospace
\usepackage{enumitem}\usepackage{amsfonts}
\usepackage{amsmath} % Recomendado para o comando \text{}
\usepackage{amssymb}
\usepackage{tikz}
\usetikzlibrary{positioning, shapes.geometric, arrows.meta}
\usepackage{listings}
\usepackage{multicol}

\lstset{
    language=Java,
    keywordstyle=\bfseries,
    showstringspaces=false,
    breaklines=true,
    frame=none,
    numbers=none,
    tabsize=2,
    commentstyle=\itshape\color{gray},
    morekeywords={this, class, private, public, String, System, out, println, Solo, Caminhao, Patio, Termometro},
    extendedchars=true,
    inputencoding=utf8
}

% TikZ Styles for ER Diagrams
\tikzset{
    block/.style={rectangle, draw, thick, minimum width=2cm, minimum height=1cm, fill=white},
    diamond_node/.style={diamond, draw, thick, minimum width=2cm, minimum height=1cm, fill=white, aspect=2},
    attribute/.style={ellipse, draw, thin, minimum width=1.5cm, minimum height=0.6cm, fill=white, font=\scriptsize},
    line/.style={draw, thick}, 
}
% ------------------------------------------------ 
% Configuração de página
% ------------------------------------------------
\geometry{
    a4paper,
    left=2cm,
    right=2cm,
    top=2cm,
    bottom=2cm,
    headheight=15pt % Adicione isso
}% Caixa para resposta discursiva
\newcommand{\resposta}[1]{
\noindent
\framebox[\textwidth]{
\begin{minipage}{0.98\textwidth}
\vspace*{#1}\null
\end{minipage}
}
}

\pagestyle{fancy}
\fancyhf{}
\renewcommand{\headrulewidth}{0pt}

\setlength{\parindent}{0pt}

% ------------------------------------------------
% Comandos auxiliares
% ------------------------------------------------
\newcolumntype{L}[1]{>{\raggedright\arraybackslash}p{#1}}
\newcolumntype{C}[1]{>{\centering\arraybackslash}p{#1}}

% ------------------------------------------------
% Documento
% ------------------------------------------------
\begin{document}

% ------------------------------------------------
% Cabeçalho - Corrigido para evitar aninhamento problemático
% ------------------------------------------------
\begin{tabularx}{\textwidth}{@{} L{4cm} X @{}}
    \raisebox{-0.5cm}{\includegraphics[width=4cm]{Universidade_de_Rio_Verde.jpg}} &
    \hfill
    \begin{tabular}{|>{\centering\arraybackslash}m{9.5cm}|>{\raggedright\arraybackslash}p{1.8cm}|}
        \hline
        \cellcolor{black}\color{white}\rule{0pt}{1.2cm}\textbf{AVALIAÇÃO 1} 
        & 
        \scriptsize Nota \tabularnewline
        \hline
    \end{tabular} 
\end{tabularx}
\vspace{0.3cm}

% ------------------------------------------------  
% Tabela principal - Estrutura limpa
% ------------------------------------------------
\begin{tabularx}{\textwidth}{|X|L{2cm}|}
\hline
\textbf{Curso} & \textbf{Campus} \\
\textit{Engenharia de Software} & \textit{Rio Verde} \\
\hline
\multicolumn{2}{|l|}{ \textit{Desenvolvimento de Software para Dispositivos Móveis}} \\
\hline
\textbf{Nome do(a) acadêmico(a)} & \textbf{Valor} \\
& 7,0 \\
\hline
\textbf{Professor} & \textbf{Turma} \\
\textit{ Sergio Souza Novak} & \textit{} \\
\hline
\textbf{Data da Avaliação} & \textbf{Código} \\
\textit{02/10/2026} & \textit{ESW417} \\
\hline
\end{tabularx}
 
% ------------------------------------------------
% Barra de Atenção
% ------------------------------------------------
\vspace{0.2cm}

\begin{tabularx}{\textwidth}{|X|}
\hline
\cellcolor{black}\color{white}
\centering \textbf{ATENÇÃO:} Somente serão passíveis de \textbf{REVISÃO} avaliações resolvidas à \textbf{TINTA}. \tabularnewline
\hline
\end{tabularx}

\vspace{0.3cm}

\textbf{ORIENTAÇÕES PARA A RESOLUÇÃO}  

\begin{itemize}
    \item A avaliação é individual, sem consulta;
    \item A interpretação de texto faz parte da avaliação;
\end{itemize}

\vspace{0.5cm}

\begin{center}
    \rule{\textwidth}{0.4pt}
    \textbf{QUESTÕES}
\end{center}

\vspace{0.5cm} 

\textbf{Código fonte de apoio para as questões 1 até 4}

\setlength{\columnseprule}{0.4pt}
\begin{multicols}{2}
\begin{lstlisting}[basicstyle=\scriptsize\ttfamily, breaklines=true]
package com.sergiosouzanovak.corpodocenteunirv


class MainActivity : ComponentActivity() {
    override fun onCreate(savedInstanceState: Bundle?) {
        super.onCreate(savedInstanceState)
        enableEdgeToEdge()
        setContent {
            CorpoDocenteUnirvTheme {
                Scaffold(modifier = Modifier.fillMaxSize()) { innerPadding ->
                    App(
                        modifier = Modifier.padding(innerPadding)
                    )
                }
            }
        }
    }
}

@Composable
fun carregarImagemAsset(nomeArquivo: String): ImageBitmap? {
    val context = LocalContext.current

    return try {
        context.assets.open(nomeArquivo).use { inputStream ->
            BitmapFactory.decodeStream(inputStream)?.asImageBitmap()
        }
    } catch (e: Exception) {
        null
    }
}

@Composable
fun PesquisadorCard(
    nome: String,
    titulacao: String,
    lattes: String,
    email: String,
    foto: String,
    modifier: Modifier = Modifier
) {
    val imagem = carregarImagemAsset(foto)

    Card(
        modifier = modifier
            .fillMaxWidth()
            .padding(16.dp),
        shape = RoundedCornerShape(4.dp),
        colors = CardDefaults.cardColors(
            containerColor = Color(0xFFF1F4F5)
        ),
        elevation = CardDefaults.cardElevation(
            defaultElevation = 0.dp
        )
    ) {
        Column(
            modifier = Modifier.width(1200.dp).height(300.dp),
            horizontalAlignment = Alignment.CenterHorizontally
        ) {

            imagem?.let {
                Image(
                    bitmap = it,
                    contentDescription = "Foto de $nome",
                    modifier = Modifier
                        .fillMaxWidth()
                        .height(120.dp),
                    contentScale = ContentScale.Crop,
                )
            }

            Text(
                text = nome,
                modifier = Modifier.padding(
                    top = 14.dp,
                    start = 8.dp,
                    end = 8.dp
                ),
                fontSize = 12.sp,
                fontWeight = FontWeight.Bold,
                color = Color(0xFF222222),
                textAlign = TextAlign.Center
            )

            Text(
                text = titulacao,
                modifier = Modifier.padding(top = 6.dp),
                fontSize = 12.sp,
                color = Color(0xFF666666),
                textAlign = TextAlign.Center
            )

            Text(
                text = lattes,
                modifier = Modifier
                    .padding(
                        top = 28.dp,
                        start = 16.dp,
                        end = 16.dp
                    )
                    .clickable {
                        // abrir Lattes
                    },
                fontSize = 12.sp,
                color = Color(0xFF007AFF),
                textAlign = TextAlign.Center
            )

            Text(
                text = email,
                modifier = Modifier.padding(
                    top = 14.dp,
                    bottom = 22.dp
                ),
                fontSize = 12.sp,
                color = Color(0xFF444444),
                textAlign = TextAlign.Center
            )
        }
    }
}

data class Docente(
    val nome: String,
    val titulacao: String,
    val lattes: String,
    val email: String,
    val foto: String
)

@Composable
fun ListaPesquisadores(){
    val docentes = listOf(
        Docente(
            nome = "ALEX MIZAEL MARTINS",
            titulacao = "Doutor",
            lattes = "www.google.com",
            email = "sergio.novak@unirv.edu.br",
            foto = "docentes/ALEX MIZAEL MARTINS.png"
        ),
        Docente(
            nome = "BACUS DE OLIVEIRA NAHIME",
            titulacao = "Doutor",
            lattes = "www.google.com",
            email = "joao@unirv.edu.br",
            foto = "docentes/BACUS DE OLIVEIRA NAHIME.png"
        ),
        Docente(
            nome = "Barbara Gomes Martins",
            titulacao = "Mestre",
            lattes = "www.google.com",
            email = "maria@unirv.edu.br",
            foto = "docentes/Barbara Gomes Martins.png"
        ),
        Docente(
            nome = "BIANCA SIMONELI DE OLIVEIRA",
            titulacao = "Doutor",
            lattes = "www.google.com",
            email = "pedro@unirv.edu.br",
            foto = "docentes/BIANCA SIMONELI DE OLIVEIRA.png"
        )
    )

    LazyVerticalGrid(
        columns = GridCells.Fixed(2),
        modifier = Modifier.fillMaxSize(),
        contentPadding = PaddingValues(16.dp),
        horizontalArrangement = Arrangement.spacedBy(16.dp),
        verticalArrangement = Arrangement.spacedBy(16.dp)
    ) {
        items(docentes) { docente ->
            PesquisadorCard(
                nome = docente.nome,
                titulacao = docente.titulacao,
                lattes = docente.lattes,
                email = docente.email,
                foto = docente.foto
            )
        }
    }
}

@Composable
fun PageTitle(){
    Column(
        horizontalAlignment = Alignment.CenterHorizontally
    ) {
        Row{
            Text(text = "Corpo Docente do curso de")
            Text(
                text = " Engenharia Civil",
                color = Color(0xFF064F9F)
            )
        }
        Row{
            Column {
                repeat(2){
                    HorizontalDivider(
                        modifier = Modifier.width(125.dp),
                        thickness = 2.dp,
                        color = Color.Gray
                    )
                    Spacer(modifier = Modifier.height(4.dp))
                }
            }
        }
        Row{
            ListaPesquisadores()
        }
    }
}

@Composable
fun App(modifier:Modifier) {
    Column{
        PageTitle()
    }
}

@Preview(showBackground = true)
@Composable
fun GreetingPreview() {
    CorpoDocenteUnirvTheme {
        App(modifier = Modifier)
    }
}
\end{lstlisting}
\end{multicols}

\vspace{0.5cm}

\textbf{Questão 1 - } No Jetpack Compose, a anotação \texttt{@Composable} é fundamental para a construção de interfaces de usuário. Qual é a principal função dessa anotação (utilizada, por exemplo, nas funções \texttt{PesquisadorCard} e \texttt{ListaPesquisadores}) no código fornecido?

\begin{enumerate}[label=\alph*)]
    \item Indicar ao compilador que a função retorna uma visualização tradicional baseada em arquivos XML do Android.
    \item Transformar a função em um componente de interface gráfica que o framework do Compose pode processar para desenhar a tela.
    \item Configurar a comunicação direta entre o aplicativo e o banco de dados interno do dispositivo móvel.
    \item Declarar que a função correspondente deve ser executada de forma assíncrona, visando otimizar a rede.
    \item Gerenciar o ciclo de vida da Activity principal (MainActivity) substituindo automaticamente o método \texttt{onCreate}.
\end{enumerate}

\vspace{0.5cm}

\textbf{Questão 2 - } Sobre os componentes de layout \texttt{Row} e \texttt{Column}, amplamente utilizados no Jetpack Compose (como visto na função \texttt{PageTitle}), julgue as afirmativas abaixo como Verdadeiras (V) ou Falsas (F):

\begin{itemize}[label={(\ \ \ )}]
    \item O componente \texttt{Column} organiza seus elementos filhos sequencialmente na vertical (um abaixo do outro).
    \item O componente \texttt{Row} organiza seus elementos filhos sequencialmente na horizontal (um ao lado do outro).
    \item Em uma \texttt{Column}, para alinhar os elementos filhos no eixo horizontal (como esquerda, centro ou direita), utiliza-se o parâmetro \texttt{horizontalAlignment}.
\end{itemize}

\vspace{0.5cm}

\textbf{Questão 3 - } No desenvolvimento Android, observa-se o uso constante de unidades como \texttt{dp} (Density-independent Pixels) para as dimensões dos componentes (ex: \texttt{16.dp}) e \texttt{sp} (Scale-independent Pixels) para o tamanho das fontes (ex: \texttt{12.sp}). Explique qual é o problema de se utilizar a medida absoluta de pixels (\texttt{px}) na construção de telas e quais são as vantagens de se adotar o \texttt{dp} e o \texttt{sp}.

\vspace{0.3cm}
\resposta{7cm}

\vspace{0.5cm}

\textbf{Questão 4 - } No código de apoio, observamos duas estruturas responsáveis por exibir componentes de interface: a função \texttt{GreetingPreview()} (anotada com \texttt{@Preview}) e a classe \texttt{MainActivity} (que herda de \texttt{ComponentActivity}). Explique a principal diferença entre os dois trechos de código. Em sua resposta, indique claramente qual deles é utilizado apenas durante a fase de desenvolvimento (para pré-visualização na IDE) e qual deles é efetivamente invocado pela aplicação real quando executada no dispositivo móvel.

\vspace{0.3cm}
\resposta{6cm}

\vspace{0.5cm}

\textbf{Questão 5 - } Uma equipe de desenvolvimento deseja criar um aplicativo para smartphones capaz de auxiliar na detecção de metais utilizando os sensores já disponíveis no dispositivo móvel. Durante o projeto, foram identificados os seguintes sensores:

\begin{itemize}[leftmargin=1.5em]
  \item Um sensor que mede a intensidade do campo magnético ao redor do celular.
  \item Um sensor que detecta movimentos e acelerações do aparelho.
  \item Um sensor que identifica a orientação e a rotação do dispositivo.
\end{itemize}

\vspace{0.4cm}

% ── Wireframes das telas do aplicativo ──────────────────────────────────────
\begin{center}
\begin{tikzpicture}[
  font=\sffamily\scriptsize,
  every node/.style={inner sep=0pt},
  phone/.style={draw=black!70, rounded corners=4mm, fill=white,
                line width=1.2pt, minimum width=3.8cm, minimum height=7cm},
  topbar/.style={draw=black!60, fill=black!10, line width=0.8pt,
                 minimum width=3.4cm, minimum height=0.55cm, rounded corners=1mm},
  widget/.style={draw=black!50, fill=white, line width=0.6pt,
                 minimum width=3.2cm, rounded corners=1mm},
  btn/.style={draw=black!70, fill=black!15, line width=0.8pt,
              minimum width=2.6cm, minimum height=0.55cm, rounded corners=2mm},
  sensorlbl/.style={draw=black!40, fill=black!5, line width=0.5pt,
                    minimum width=3.2cm, minimum height=0.45cm, rounded corners=1mm},
]

%% ─────── TELA 1: Escaneamento ───────
\node[phone] (phone1) at (0,0) {};

% câmera (notch)
\draw[fill=black!25, draw=black!50, line width=0.5pt]
  (0, 3.3) circle (0.12cm);

% barra de status
\node[topbar, anchor=north] (bar1) at (0, 3.1)
  {\textbf{MetalScan} \hfill 09:41};

% ícone de campo magnético (ondas)
\begin{scope}[shift={(0, 1.8)}]
  \draw[line width=1.2pt] (0,0) circle (0.38cm);
  \draw[line width=0.7pt] (0,0) circle (0.62cm);
  \draw[line width=0.4pt, dashed] (0,0) circle (0.86cm);
  \node[font=\sffamily\tiny\bfseries] at (0,-1.1) {Campo Magnético};
\end{scope}

% barra de intensidade
\node[widget, minimum height=0.38cm, anchor=north] (bar_bg) at (0, 0.55)
  {};
\node[draw=none, fill=black!30, minimum width=1.1cm, minimum height=0.28cm,
      anchor=west, rounded corners=0.5mm] at ([xshift=0.05cm]bar_bg.west) {};
\node[font=\sffamily\tiny, anchor=east] at ([xshift=-0.08cm]bar_bg.east)
  {intensidade};

% leituras dos sensores
\node[sensorlbl, anchor=north] (s1) at (0, 0.1)
  {\textbf{Magnetômetro:} \ 48 µT};
\node[sensorlbl, anchor=north] (s2b) at (0, -0.45)
  {\textbf{Acelerômetro:} \ 0.3 m/s²};
\node[sensorlbl, anchor=north] (s3) at (0, -0.98)
  {\textbf{Giroscópio:} \ 2.1 °/s};

% botão iniciar
\node[btn, anchor=north] at (0, -1.6)
  {\textbf{\ Iniciar Varredura}};

% legenda
\node[font=\sffamily\tiny\itshape, text=black!60, anchor=north] at (0,-3.0)
  {Tela 1 — Escaneamento};

%% ─────── TELA 2: Metal Detectado ───────
\node[phone] (phone2) at (5.5,0) {};

% câmera
\draw[fill=black!25, draw=black!50, line width=0.5pt]
  (5.5, 3.3) circle (0.12cm);

% barra de status
\node[topbar, anchor=north] (bar2) at (5.5, 3.1)
  {\textbf{MetalScan} \hfill 09:42};

% ícone de alerta
\begin{scope}[shift={(5.5, 1.75)}]
  \draw[line width=1.5pt, fill=black!8]
    (0, 0.55) -- (-0.55,-0.45) -- (0.55,-0.45) -- cycle;
  \node[font=\bfseries\small] at (0, 0.04) {!};
  \node[font=\sffamily\tiny\bfseries] at (0,-0.85)
    {METAL DETECTADO};
\end{scope}

% barra de intensidade alta
\node[widget, minimum height=0.38cm, anchor=north] (bar_bg2) at (5.5, 0.55) {};
\node[draw=none, fill=black!60, minimum width=2.8cm, minimum height=0.28cm,
      anchor=west, rounded corners=0.5mm]
  at ([xshift=0.05cm]bar_bg2.west) {};
\node[font=\sffamily\tiny, anchor=east, text=white]
  at ([xshift=-0.12cm]bar_bg2.east) {alta};

% leituras — estado de alerta
\node[sensorlbl, anchor=north] at (5.5, 0.1)
  {\textbf{Magnetômetro:} \ 312 µT};
\node[sensorlbl, anchor=north] at (5.5, -0.45)
  {\textbf{Acelerômetro:} \ 1.9 m/s²};
\node[sensorlbl, anchor=north] at (5.5, -0.98)
  {\textbf{Giroscópio:} \ 8.4 °/s};

% botão nova varredura
\node[btn, anchor=north] at (5.5, -1.6)
  {\textbf{\ Nova Varredura}};

% legenda
\node[font=\sffamily\tiny\itshape, text=black!60, anchor=north] at (5.5,-3.0)
  {Tela 2 — Metal Detectado};
 
%% seta de transição
\draw[-{Stealth[length=6pt]}, line width=1pt, black!50]
  ([xshift=0.25cm]phone1.east) -- ([xshift=-0.25cm]phone2.west)
  node[midway, above, font=\sffamily\tiny, text=black!50] {detecção};

\end{tikzpicture}
\end{center}
% ────────────────────────────────────────────────────────────────────────────

\vspace{0.3cm}
Considerando o funcionamento desse aplicativo para detecção de metais, assinale a alternativa que indica corretamente qual é o principal sensor responsável por medir as variações e identificar a presença do metal ao redor do aparelho:

\vspace{0.4cm}

\begin{enumerate}[label=\alph*)]
    \item Acelerômetro
    \item Magnetômetro
    \item Giroscópio
    \item Sensor de Proximidade
    \item Barômetro
\end{enumerate}

\vspace{0.5cm}

\textbf{Questão 6 - } Analise o seguinte trecho de código escrito em Kotlin:

\begin{lstlisting}[basicstyle=\small\ttfamily, breaklines=true]
val soma = { a: Int, b: Int -> a + b }

fun main() {
    println(soma(5, 3)) // Saida: 8
}
\end{lstlisting}

O trecho de código acima demonstra uma importante característica da linguagem Kotlin. Assinale a alternativa que descreve corretamente o conceito utilizado na declaração da variável \texttt{soma}:

\begin{enumerate}[label=\alph*)]
    \item Tratamento de exceções utilizando blocos anônimos para capturar erros matemáticos.
    \item Declaração de uma função de extensão estática na classe \texttt{Int}.
    \item Instanciação de uma classe abstrata responsável pelas operações aritméticas básicas.
    \item Definição de uma função lambda (função anônima), que é tratada como um tipo de primeira classe e pode ser armazenada em variáveis.
    \item Sobrecarga do operador de adição nativo da linguagem para modificar seu funcionamento.
\end{enumerate}

\vspace{0.5cm}

\textbf{Questão 7 - } Observe o trecho de código Kotlin abaixo, que apresenta diferentes formas de inicializar coleções de dados:

\begin{lstlisting}[basicstyle=\small\ttfamily, breaklines=true]
val lista = mutableListOf("Item 1")
val lista2 = arrayListOf("Item 1") // nativo

// Set (HashSet)
val conjunto = mutableSetOf(10)
val conjunto2 = hashSetOf(10)

// Map (HashMap)
val mapa = mutableMapOf("Chave" to 100)
val mapa2 = hashMapOf("Chave" to 100)
\end{lstlisting}

A linguagem Kotlin provê diversas funções utilitárias para a criação de coleções. Com base nas estruturas apresentadas (\texttt{List}, \texttt{Set} e \texttt{Map}) e nos conceitos de coleções em Kotlin, assinale a afirmativa correta:

\begin{enumerate}[label=\alph*)]
    \item As funções que possuem o prefixo \texttt{mutable} criam coleções de leitura restrita (imutáveis), cujos elementos não podem ser alterados após a instanciação.
    \item A estrutura \texttt{Set} (como a criada por \texttt{mutableSetOf} ou \texttt{hashSetOf}) é projetada para permitir o armazenamento de vários elementos idênticos e duplicados.
    \item A estrutura \texttt{Map} armazena elementos utilizando o paradigma de associação chave-valor, onde cada chave mapeada deve ser necessariamente única na coleção.
    \item O uso da palavra-chave \texttt{val} na declaração das variáveis impede que novos elementos sejam adicionados ou removidos das listas e conjuntos mutáveis.
    \item Em Kotlin, não existe diferença conceitual entre \texttt{List} e \texttt{Set}, pois ambas utilizam índices inteiros, não possuem ordem e não permitem repetidos.
\end{enumerate}

\vspace{0.5cm}

\textbf{Questão 8 - } A linguagem Kotlin oferece diversas funções de alta ordem (higher-order functions) que facilitam a manipulação e transformação de coleções de forma funcional, sendo \texttt{filter} e \texttt{map} algumas das mais utilizadas. Sobre o tema, julgue as afirmativas abaixo como Verdadeiras (V) ou Falsas (F):

\begin{itemize}[label={(\ \ \ )}]
    \item A função \texttt{filter} é utilizada para percorrer uma coleção e retornar uma nova coleção contendo apenas os elementos que satisfazem uma dada condição (predicado).
    \item A função \texttt{map} tem o objetivo de transformar os itens de uma coleção, aplicando uma função fornecida a cada elemento original e retornando uma nova coleção com os resultados gerados.
    \item O uso das funções \texttt{filter} e \texttt{map} altera e substitui diretamente os valores dentro da coleção original na memória, não sendo necessário armazenar o retorno.
    \item Ao utilizar expressões lambdas no Kotlin (como dentro do bloco do \texttt{filter} ou \texttt{map}), o identificador implícito \texttt{it} pode ser utilizado para referenciar o elemento atual da iteração, encurtando o código.
\end{itemize}

\end{document}